\documentclass[a4paper]{article}
% Shared preamble for the labs (compile from Labs/ with XeLaTeX or tectonic).
% Same fonts and header as the Series sheets, plus colored boxes and tables.

\usepackage{enumitem}
\usepackage[margin=2.2cm]{geometry}
\usepackage{amsmath,amsfonts,amssymb}
\usepackage[charter]{mathdesign}
\usepackage{array}
\usepackage{colortbl}
\usepackage{fancyhdr}
\usepackage[skins,breakable]{tcolorbox}
\usepackage{polyglossia}

%------------------------------------------------
% Languages and fonts
%------------------------------------------------
\setdefaultlanguage[calendar=gregorian, numerals=maghrib]{arabic}
\setotherlanguage{english}

\IfFontExistsTF{Amiri}{%
    \newfontfamily\arabicfont[Script=Arabic, Scale=1.15]{Amiri}
}{%
    \newfontfamily\arabicfont[
        Script=Arabic, Scale=1.15,
        Path=../fonts/, Extension=.ttf,
        UprightFont=*-Regular, BoldFont=*-Bold,
        ItalicFont=*-Italic, BoldItalicFont=*-BoldItalic
    ]{Amiri}
}

\IfFontExistsTF{AlBattar}{%
    \newfontfamily\arabicfontsf[Script=Arabic, Scale=1.2]{AlBattar}
    \newfontfamily\khatt[Script=Arabic, Scale=1.2]{AlBattar}
}{%
    \newfontfamily\arabicfontsf[
        Script=Arabic, Scale=1.2, Path=../fonts/,
        BoldFont=ae_AlBattar.ttf
    ]{ae_AlBattar.ttf}
    \newfontfamily\khatt[
        Script=Arabic, Scale=1.2, Path=../fonts/,
        BoldFont=ae_AlBattar.ttf
    ]{ae_AlBattar.ttf}
}

%------------------------------------------------
% Colors (same palette as the slides)
%------------------------------------------------
\definecolor{brandTeal}{HTML}{00B4D8}
\definecolor{brandOrange}{HTML}{FF7043}
\definecolor{brandDark}{HTML}{263238}
\definecolor{brandMuted}{HTML}{6C757D}
\colorlet{tealDeep}{brandTeal!70!black}

%------------------------------------------------
% Page setup
%------------------------------------------------
\renewcommand{\footrulewidth}{1pt}
\renewcommand{\headrulewidth}{1pt}
\rhead{\khatt المدرسة العليا للأساتذة ببوزريعة}
\lhead{\khatt ملحقة بومرداس}
\cfoot{\thepage}
\lfoot{\khatt الأستاذ: د. بن عبد الرزاق يوسف}
\parindent 0pt
\setlength{\headsep}{15pt}
\setlength{\headheight}{16pt}
\setlength{\parskip}{4pt}
\pagestyle{fancy}

\setlist[enumerate,1]{label=\arabic*), itemsep=4pt, topsep=3pt}
\setlist[enumerate,2]{label=\alph*), itemsep=2pt, topsep=2pt}
\setlist[itemize]{itemsep=2pt, topsep=2pt}

%------------------------------------------------
% Tables
%------------------------------------------------
\renewcommand{\arraystretch}{1.4}
\setlength{\arrayrulewidth}{0.6pt}
\arrayrulecolor{brandMuted!80}
% Header cell
\newcommand{\ras}[1]{\cellcolor{brandTeal!20}\textbf{#1}}
% Highlighted result cell
\newcommand{\natija}[1]{\cellcolor{brandOrange!15}\textbf{#1}}
% Empty cell to be filled by the student
\newcommand{\fragh}{\rule{0pt}{0.55cm}}
\newcolumntype{C}[1]{>{\centering\arraybackslash}m{#1}}
\newcolumntype{R}[1]{>{\raggedleft\arraybackslash}m{#1}}

%------------------------------------------------
% Boxes
%------------------------------------------------
\tcbset{labbox/.style={enhanced, breakable, arc=3pt, boxrule=1.2pt,
    coltitle=white, fonttitle=\khatt\bfseries,
    top=1.5mm, bottom=1.5mm, left=3mm, right=3mm,
    before skip=6pt, after skip=6pt}}
% Worked example
\newtcolorbox{mithal}[1]{labbox, colback=brandTeal!6, colframe=brandTeal,
    title={مثال محلول: #1}}
% Reminder / method
\newtcolorbox{tadhkir}[1][تذكير]{labbox, colback=brandTeal!3,
    colframe=tealDeep, title={#1}}
% Warning
\newtcolorbox{tanbih}[1][انتبه]{labbox, colback=brandOrange!8,
    colframe=brandOrange, title={#1}}
% Exercise
\newtcolorbox{tamrin}[1]{labbox, colback=white, colframe=brandDark!80,
    title={#1}}

% Section title
\newcommand{\qism}[1]{\section*{\khatt\underline{#1}}}

% Title block: \unwan{lab number}{title}{subtitle}
\newcommand{\unwan}[3]{%
    \begin{center}
        {\Large\khatt\textbf{المخبر التطبيقي رقم #1}}\\[0.15cm]
        {\Large\khatt\textbf{#2}}\\[0.1cm]
        {\khatt\textbf{#3}}
    \end{center}
    \rule{\textwidth}{1pt}
    \begin{center}
        \small المدة: ساعة ونصف \quad$\cdot$\quad الأدوات: آلة حاسبة، أداة ذكاء اصطناعي، مصدر موثوق
    \end{center}
    \rule{\textwidth}{1pt}}

% Percentage reminder used in every lab
\newcommand{\nisba}{%
    \begin{tadhkir}[تذكير: حساب النسبة المئوية]
        \begin{center}
            \textbf{النسبة المئوية = (الجزء $\div$ الكل) $\times$ 100}
        \end{center}
        مثال: 15 إجابة صحيحة من 20 سؤالًا:
        $\dfrac{15}{20}\times 100 = 75\%$.
    \end{tadhkir}}


\newcommand{\yes}{\textcolor{tealDeep}{$\checkmark$}}
\newcommand{\no}{\textcolor{brandOrange}{$\times$}}

\begin{document}

\unwan{3}{التفاعل الفعّال مع الذكاء الاصطناعي}{تقييم الأوامر بالأرقام ــ شبكة العناصر الخمسة}

\qism{أهداف المخبر}
\begin{itemize}
    \item تقييم أي أمر بشبكة العناصر الخمسة (نقطة لكل عنصر، من 5).
    \item ربط نقطة الأمر بجودة الإجابة.
    \item قياس احترام الأداة للقيود (عدد الكلمات، عدد الأمثلة).
    \item تحليل نتائج استبيان حول أوامر الطلبة.
\end{itemize}

\begin{tadhkir}[شبكة التقييم: العناصر الخمسة]
    نعطي \textbf{نقطة واحدة} لكل عنصر موجود في الأمر، فتكون النقطة القصوى \textbf{5}:
    \begin{center}
    \begin{tabular}{|C{2.4cm}|C{2.4cm}|C{2.4cm}|C{2.4cm}|C{2.4cm}|}
        \hline
        \ras{الدور} & \ras{المهمة} & \ras{السياق} & \ras{القيود} & \ras{التنسيق} \\ \hline
        من تكون الأداة؟ & ماذا تفعل؟ & لمن؟ وفي أي ظرف؟ & ما الشروط؟ & ما شكل الإجابة؟ \\ \hline
    \end{tabular}
    \end{center}
\end{tadhkir}

\nisba

%================================================
\qism{الجزء الأول: أمثلة محلولة}
%================================================

\begin{mithal}{تقييم ثلاثة أوامر}
\textbf{الأمر 1:} «اشرح الحال.»\\
\textbf{الأمر 2:} «اشرح الحال لتلاميذ السنة الثانية متوسط في 150 كلمة.»\\
\textbf{الأمر 3:} «تصرّف كأستاذ لغة عربية. اشرح الحال لتلاميذ السنة الثانية متوسط،
في 150 كلمة، ونظّم الإجابة في ثلاثة أجزاء: تعريف، أمثلة، خلاصة.»

\begin{center}
\begin{tabular}{|C{1.2cm}|C{1.3cm}|C{1.3cm}|C{1.3cm}|C{1.3cm}|C{1.3cm}|C{1.6cm}|C{2.2cm}|}
    \hline
    \ras{الأمر} & \ras{الدور} & \ras{المهمة} & \ras{السياق} & \ras{القيود} & \ras{التنسيق}
        & \ras{النقطة /5} & \ras{جودة الإجابة /10} \\ \hline
    1 & \no  & \yes & \no  & \no  & \no  & 1 & 3 \\ \hline
    2 & \no  & \yes & \yes & \yes & \no  & 3 & 6 \\ \hline
    3 & \yes & \yes & \yes & \yes & \yes & \natija{5} & \natija{9} \\ \hline
\end{tabular}
\end{center}

\textbf{الحساب:}
\begin{itemize}
    \item متوسط نقاط الأوامر الثلاثة: $\dfrac{1+3+5}{3} = 3$ من 5.
    \item من الأمر 1 إلى الأمر 3: ربحنا $5-1 = 4$ عناصر، وارتفعت الجودة $9-3=6$ نقاط.
\end{itemize}
\textbf{الاستنتاج:} كلما زاد عدد العناصر في الأمر، تحسنت جودة الإجابة.
\end{mithal}

\begin{mithal}{أي عنصر ينساه الطلبة؟}
جمع أستاذ \textbf{30 أمرًا} كتبها طلبته، وعدّ كم أمرًا يحتوي على كل عنصر:

\begin{center}
\begin{tabular}{|R{3cm}|C{3cm}|C{3cm}|}
    \hline
    \ras{العنصر} & \ras{عدد الأوامر} & \ras{النسبة} \\ \hline
    المهمة   & 30 & 100\% \\ \hline
    السياق   & 15 & 50\%  \\ \hline
    القيود   & 12 & 40\%  \\ \hline
    الدور    & 9  & 30\%  \\ \hline
    التنسيق  & 6  & \natija{20\%} \\ \hline
\end{tabular}
\end{center}

\textbf{الحساب:} مثلًا التنسيق: $\dfrac{6}{30}\times 100 = 20\%$.\\
\textbf{الاستنتاج:} كل الطلبة يكتبون المهمة، لكن \textbf{التنسيق} هو العنصر الأكثر نسيانًا
(4 أوامر من كل 5 لا تحدد شكل الإجابة). يجب أن يركز الأستاذ عليه في الحصة القادمة.
\end{mithal}

%================================================
\qism{الجزء الثاني: تمارين تطبيقية}
%================================================

\begin{tamrin}{التمرين 1: قيّم الأوامر}
ضع \yes\ أو \no\ في كل خانة، ثم احسب النقطة:

\begin{enumerate}
    \item «لخّص هذه القصيدة.»
    \item «تصرّف كأستاذ بلاغة واشرح الاستعارة.»
    \item «اقترح عشر جمل فيها مفعول لأجله، لتلاميذ السنة الثالثة متوسط.»
    \item «أنشئ خمسة أسئلة فهم حول النص الآتي، مع الإجابات، في جدول من عمودين.»
    \item «أنا أستاذ متربص وتلاميذي في السنة الأولى متوسط. تصرّف كمفتش تربوي،
          واقترح خطة درس للنعت مدتها 45 دقيقة، في شكل جدول.»
    \item «اكتب شيئًا جميلًا.»
\end{enumerate}

\begin{center}
\begin{tabular}{|C{1.4cm}|C{1.6cm}|C{1.6cm}|C{1.6cm}|C{1.6cm}|C{1.6cm}|C{1.8cm}|}
    \hline
    \ras{الأمر} & \ras{الدور} & \ras{المهمة} & \ras{السياق} & \ras{القيود} & \ras{التنسيق} & \ras{النقطة /5} \\ \hline
    1 & \fragh & & & & & \\ \hline
    2 & \fragh & & & & & \\ \hline
    3 & \fragh & & & & & \\ \hline
    4 & \fragh & & & & & \\ \hline
    5 & \fragh & & & & & \\ \hline
    6 & \fragh & & & & & \\ \hline
\end{tabular}
\end{center}

\begin{enumerate}[resume]
    \item احسب متوسط نقاط الأوامر الستة.
    \item أعد كتابة الأمر 1 ليحصل على 5 من 5.
\end{enumerate}
\end{tamrin}

\begin{tamrin}{التمرين 2: هل احترمت الأداة القيد؟}
طلب أستاذ شرحًا لدرس التمييز \textbf{في 150 كلمة}، وجرّب ثلاث صياغات:

\begin{center}
\begin{tabular}{|C{1.8cm}|C{2.6cm}|C{2.6cm}|C{2.8cm}|C{2.6cm}|}
    \hline
    \ras{المحاولة} & \ras{عدد الكلمات} & \ras{الزيادة (كلمة)} & \ras{نسبة الزيادة} & \ras{مقبولة؟} \\ \hline
    1 & 230 & \fragh & & \\ \hline
    2 & 180 & \fragh & & \\ \hline
    3 & 155 & \fragh & & \\ \hline
\end{tabular}
\end{center}

\begin{enumerate}
    \item أكمل الجدول. (نسبة الزيادة = الزيادة $\div$ 150 $\times$ 100)
    \item يقبل الأستاذ زيادة لا تتجاوز \textbf{10\%}. ما المحاولات المقبولة؟
    \item في المحاولة 3 أضاف الأستاذ عبارة «لا تتجاوز 150 كلمة بأي حال». ماذا تستنتج؟
\end{enumerate}
\end{tamrin}

\begin{tamrin}{التمرين 3: التحسين خطوة بخطوة}
بنى طالب أمره على أربع درجات، وقيّم كل إجابة من 10:

\begin{center}
\begin{tabular}{|C{1.6cm}|R{4cm}|C{2.4cm}|C{2.4cm}|C{2.6cm}|}
    \hline
    \ras{الدرجة} & \ras{ما أضافه} & \ras{التقييم /10} & \ras{التحسن} & \ras{الزمن (دقيقة)} \\ \hline
    1 & أمر عام              & 3 & -- & 1 \\ \hline
    2 & + السياق             & 5 & \fragh & 2 \\ \hline
    3 & + المهمة والقيود     & 7 & \fragh & 3 \\ \hline
    4 & + التنسيق            & 9 & \fragh & 3 \\ \hline
\end{tabular}
\end{center}

\begin{enumerate}
    \item أكمل عمود التحسن (الفرق مع الدرجة السابقة).
    \item ما الزمن الكلي؟ وما مجموع التحسن من الدرجة 1 إلى الدرجة 4؟
    \item هل تستحق 9 دقائق هذا التحسن؟ علّل.
    \item لو حصل الطالب في الدرجة 5 على 9 من 10 أيضًا، هل يواصل التعديل؟ لماذا؟
\end{enumerate}
\end{tamrin}

\begin{tamrin}{التمرين 4: استبيان في قسمك}
حلّل أستاذ \textbf{40 أمرًا} كتبها طلبة السنة الأولى:

\begin{center}
\begin{tabular}{|R{3cm}|C{3cm}|C{3cm}|}
    \hline
    \ras{العنصر} & \ras{عدد الأوامر} & \ras{النسبة} \\ \hline
    الدور    & 10 & \fragh \\ \hline
    المهمة   & 40 & \fragh \\ \hline
    السياق   & 22 & \fragh \\ \hline
    القيود   & 18 & \fragh \\ \hline
    التنسيق  & 8  & \fragh \\ \hline
\end{tabular}
\end{center}

\begin{enumerate}
    \item أكمل عمود النسبة.
    \item رتّب العناصر من الأكثر حضورًا إلى الأقل.
    \item ما العنصران اللذان يجب التدريب عليهما أولًا؟
    \item \textbf{تطبيق:} اجمع أوامر 10 زملاء في قسمك، وأنشئ جدولًا مماثلًا.
\end{enumerate}
\end{tamrin}

\begin{tamrin}{التمرين 5: تجربة على الأداة}
اختر درسًا (مثل: الاستعارة، المفعول المطلق، كان وأخواتها).
اكتب أمره في أربع درجات كما في الدرس، ونفّذ كل درجة، ثم أكمل الجدول:

\begin{center}
\begin{tabular}{|C{1.8cm}|C{2.4cm}|C{2.6cm}|C{2.6cm}|C{2.6cm}|}
    \hline
    \ras{الدرجة} & \ras{النقطة /5} & \ras{عدد الكلمات} & \ras{عدد الأمثلة} & \ras{تقييمك /10} \\ \hline
    1 & \fragh & & & \\ \hline
    2 & \fragh & & & \\ \hline
    3 & \fragh & & & \\ \hline
    4 & \fragh & & & \\ \hline
\end{tabular}
\end{center}
هل تؤكد نتائجك العلاقة بين نقطة الأمر وجودة الإجابة؟
\end{tamrin}

%================================================
\newpage
\qism{الحلول المختصرة}
%================================================

\textbf{التمرين 1:}
\begin{center}
\begin{tabular}{|C{1.4cm}|C{1.6cm}|C{1.6cm}|C{1.6cm}|C{1.6cm}|C{1.6cm}|C{1.8cm}|}
    \hline
    \ras{الأمر} & \ras{الدور} & \ras{المهمة} & \ras{السياق} & \ras{القيود} & \ras{التنسيق} & \ras{النقطة /5} \\ \hline
    1 & \no  & \yes & \no  & \no  & \no  & 1 \\ \hline
    2 & \yes & \yes & \no  & \no  & \no  & 2 \\ \hline
    3 & \no  & \yes & \yes & \yes & \no  & 3 \\ \hline
    4 & \no  & \yes & \no  & \yes & \yes & 3 \\ \hline
    5 & \yes & \yes & \yes & \yes & \yes & 5 \\ \hline
    6 & \no  & \yes & \no  & \no  & \no  & 1 \\ \hline
\end{tabular}
\end{center}
(في الأمر 6 المهمة موجودة لكنها غامضة جدًا.)
المتوسط: $\dfrac{1+2+3+3+5+1}{6} = \dfrac{15}{6} = 2.5$ من 5.\\
مثال لإعادة كتابة الأمر 1: «تصرّف كأستاذ أدب عربي. لخّص هذه القصيدة لطلبة السنة الأولى
جامعي، في 100 كلمة، ثم اذكر الفكرة العامة في سطر واحد.»

\textbf{التمرين 2:}
الزيادة: 80، 30، 5 كلمات. النسبة: $53.3\%$، $20\%$، $3.3\%$.
المقبولة: المحاولة 3 فقط.
القيد الواضح والصريح يجعل الأداة تحترم الطول المطلوب أكثر.

\textbf{التمرين 3:}
التحسن: $+2$، $+2$، $+2$. الزمن الكلي: $1+2+3+3 = 9$ دقائق. مجموع التحسن: $9-3 = 6$ نقاط.
نعم، ربح 6 نقاط في 9 دقائق مفيد. لا يواصل إذا لم تتحسن النتيجة: الإجابة صارت كافية.

\textbf{التمرين 4:}
النسب: الدور $25\%$، المهمة $100\%$، السياق $55\%$، القيود $45\%$، التنسيق $20\%$.
الترتيب: المهمة، السياق، القيود، الدور، التنسيق.
التدريب أولًا على: \textbf{التنسيق} و\textbf{الدور}.

\textbf{التمرين 5:} إجابة شخصية مبنية على قياسات حقيقية.

\end{document}
